\section{Discussion}
\label{sec:discussion}

\subsection{Why direct SOM-QE refinement did not outperform Gradient}

The controlled adaptive-mesh experiments provide a clear result for the
present lid-driven-cavity problem: direct refinement based on SOM
quantization error does not outperform the conventional velocity-gradient
indicator.

This result is particularly relevant because the preliminary indicator
analysis initially suggested that SOM-QE contained potentially useful
information. On the M40 pilot solution, quantization error was spatially
correlated with regions exhibiting relatively large discrepancies with the
M160 reference. At some refinement budgets, ranking cells by SOM-QE captured
a substantial fraction of the initial M40--M160 discrepancy.

However, the adaptive simulations demonstrate that identifying cells with
large current discrepancy is not equivalent to identifying cells whose
refinement will produce the largest reduction in the final global error.

The distinction may be written conceptually as

\begin{equation}
\text{current local discrepancy}
\neq
\text{refinement utility}.
\end{equation}

Refinement modifies the discrete CFD system itself. The resulting velocity
field is globally coupled through the governing equations, pressure--velocity
interaction, boundary conditions, and neighboring control volumes.
Consequently, refining one region can improve the solution in locations that
were not directly selected for refinement.

This behavior was observed explicitly in the 10\% refinement-benefit
analysis. Gradient refinement improved the solution over a broad fraction of
the common diagnostic grid, including regions not directly selected by the
Gradient criterion. SOM-QE refinement produced a smaller global benefit and
even increased the local discrepancy in portions of the unselected domain.

Therefore, the value of a refinement indicator cannot be established solely
from its correlation with an existing error field. It must ultimately be
evaluated through the solution obtained after refinement.

\subsection{What the SOM actually learned}

The inferior performance of direct SOM-QE refinement should not be
interpreted as evidence that the SOM failed to learn meaningful information
about the flow.

The SOM was trained using local flow variables and velocity-gradient
information without including the spatial coordinates \(x/L\) and \(y/L\)
as input features. Nevertheless, the resulting best-matching-unit
classification formed spatially coherent regions when mapped back onto the
cavity.

This indicates that the SOM identified recurring local flow states rather
than merely memorizing spatial position.

Analysis of the SOM neurons further showed that cells having comparable
velocity-gradient magnitudes could belong to different neurons because their
velocity magnitude and directional components differed. In this sense, the
Gradient indicator and the SOM encode different types of information:

\begin{equation}
\text{Gradient}
\rightarrow
\text{strength of local solution variation},
\end{equation}

whereas

\begin{equation}
\text{SOM state}
\rightarrow
\text{type of local flow regime}.
\end{equation}

This distinction explains why SOM-based information can be physically
structured without automatically becoming an effective numerical-error
indicator.

\subsection{Quantization error is not an error estimator}

SOM quantization error measures the distance between an input sample and its
best-matching SOM prototype in standardized feature space. A large value
therefore indicates that the local flow state is represented relatively
poorly by the trained map.

It does not, by definition, measure discretization error.

In the present case, SOM-QE and velocity-gradient magnitude were strongly
correlated. Their cell-level correlation was approximately

\begin{equation}
r\left(|\nabla U|,QE\right)=0.879.
\end{equation}

Both quantities therefore respond strongly to extreme flow states, including
the upper-corner regions associated with the moving-lid boundary condition.

This explains why SOM-QE initially appeared promising as an error-location
indicator. Nevertheless, the adaptive experiments separate correlation from
causation: the cells with the largest quantization error are not necessarily
the cells whose refinement most efficiently reduces the global CFD
discrepancy.

For the present experiment, the hypothesis that SOM quantization error can
serve directly as a superior refinement indicator is therefore not
supported.

\subsection{Interpretation of the final adaptive meshes}

The final 10\% adaptive meshes provide a geometric interpretation of the
quantitative results.

Both SOM-QE and Gradient begin from the same M40 pilot mesh, select the same
number of cells for refinement, and produce final meshes containing 2080
cells. Nevertheless, Figure~\ref{fig:adaptive_meshes} shows that the cells
are distributed differently.

The Gradient criterion concentrates a larger fraction of the available
resolution close to the moving lid and upper cavity, where the velocity
field changes rapidly. SOM-QE allocates part of the same finite refinement
budget to additional regions farther from the lid.

Those additional SOM-selected regions are not random. They correspond to
flow states that the SOM regards as unusual or poorly represented.
The adaptive CFD results nevertheless show that resolving those states more
finely does not produce as much global numerical benefit as concentrating
the same computational resources according to the Gradient criterion.

This provides an important engineering distinction:

\begin{equation}
\text{physically interesting}
\neq
\text{numerically profitable to refine}.
\end{equation}

For adaptive CFD, the objective is not merely to identify interesting flow
structures, but to allocate a limited computational budget where additional
resolution most improves the quantities of interest.

\subsection{What was learned from Hybrid-v1}

The Hybrid-v1 indicator was deliberately designed after rejecting direct
SOM-QE refinement as the primary strategy. Instead of allowing SOM
information to replace the conventional indicator, the velocity gradient
remained dominant and the SOM regime-transition indicator only modified its
priority:

\begin{equation}
I_{\mathrm{hyb},i}
=
G_i^{*}
\left(
1+\alpha T_i
\right),
\qquad
\alpha=0.5.
\end{equation}

The coefficient was fixed before evaluating the hybrid cases against the
M160 reference. It was not optimized using the reference solution.

This formulation produced a controlled perturbation of the Gradient
selection. Depending on refinement budget, approximately 91--96\% of the
Gradient-selected cells were retained.

The cells introduced by Hybrid-v1 consistently had somewhat smaller velocity
gradients but substantially larger SOM regime-transition indicators than the
cells they replaced. The algorithm therefore behaved according to its
intended design.

At the 10\% budget, changing only 11 of the 160 Gradient-selected cells
produced an approximately 1.24\% reduction in RMS discrepancy. However, the
same mechanism did not improve RMS performance at the 5\%, 15\%, or 20\%
budgets.

The appropriate conclusion is therefore not that the hybrid method improves
Gradient refinement. Rather, the experiment demonstrates that SOM-derived
flow-state information can alter refinement decisions in a controlled and
physically interpretable way, while the numerical value of those changes
remains dependent on the specific cells exchanged.

\subsection{Refinement utility as the central problem}

The combined results suggest a more precise research question for future
work.

The original question can be formulated as whether a SOM can identify cells
that should be refined. The present results indicate that this formulation
is too broad. A more useful question is:

\begin{quote}
What information learned by a SOM, if any, is predictive of the numerical
utility of refining a cell or region?
\end{quote}

This distinction changes the role of machine learning in the adaptive
workflow.

Rather than asking an unsupervised model to replace established numerical
indicators, a more defensible strategy is to use conventional physics-based
or numerical indicators as the primary refinement signal and investigate
whether learned flow-state information can improve decisions near the
refinement threshold.

The Hybrid-v1 experiment represents a first controlled test of this concept,
but the present evidence is insufficient to establish a general advantage.


\subsection{Reframing the SOM evaluation}

The direct SOM-QE results changed the research question rather than ending
the investigation. Quantization error was shown not to be a reliable direct
proxy for discretization error or refinement utility. The subsequent
analysis therefore used the SOM in a different role: identifying interfaces
between distinct learned flow states.

This distinction is important. A SOM does not need to predict numerical
error directly for its representation of the flow to be useful. It may
instead identify spatial structures that complement a conventional
numerical indicator. The SOM-interface criterion was constructed around
this hypothesis.

The new indicator combines three pieces of information available from the
M40 pilot solution:

\begin{equation}
I_{\mathrm{SI}}
=
T\sqrt{G^*\Delta U^*}.
\end{equation}

The transition term identifies a change in learned flow regime, while the
gradient and velocity-jump terms retain sensitivity to the local strength of
the physical variation. The resulting refinement is therefore not a pure
machine-learning replacement of Gradient. It is a physics-informed use of
the SOM representation.

\subsection{Cell-matched evaluation of SOM-interface}

A direct comparison at a nominal 10\% budget would have been misleading for
the new criterion because fixed-seed graph expansion naturally produces a
different number of selected cells. The experiments were therefore followed
by cell-matched Gradient controls.

The \(1N\) SOM-interface case selected 78 cells and was compared with a
Gradient selection of exactly 78 cells. The \(2N\) case selected 114 cells
and was compared with exactly 114 Gradient-selected cells. In both cases,
the final meshes therefore contained identical numbers of cells.

Under this controlled comparison, SOM-interface produced lower
volume-weighted mean and RMS discrepancies than the corresponding Gradient
case. The mean discrepancy reduction was approximately \(13.9\%\) for the
\(1N\) comparison and \(18.7\%\) for the \(2N\) comparison. The RMS reduction
was smaller, approximately \(1.4\%\) and \(4.0\%\), respectively.

The result is noteworthy because the methods shared a substantial fraction
of their selected cells, yet allocated the remaining cells differently.
The overlap decreased from \(67.95\%\) in the \(1N\) comparison to \(64.91\%\)
in the \(2N\) comparison.

The observed improvement should nevertheless be interpreted carefully. The
maximum pointwise discrepancy remained essentially unchanged between the
paired cases, and the simulations represent only one Reynolds number and
one cavity geometry. The present evidence therefore supports the statement
that the SOM-interface hypothesis is promising for this controlled test,
not that the method has been established as a generally superior adaptive
strategy.

\subsection{Why the new result is different from SOM-QE}

The contrast between SOM-QE and SOM-interface is itself an important result.
SOM-QE ranks cells according to how poorly their local feature vector is
represented by a SOM prototype. SOM-interface instead uses the spatial
organization of the BMU classification.

Consequently,

\begin{equation}
QE
\rightarrow
\text{representation mismatch},
\end{equation}

whereas

\begin{equation}
T
\rightarrow
\text{learned flow-state transition}.
\end{equation}

The second-stage criterion further combines the transition with local
physical variation. The new result therefore does not contradict the
earlier negative SOM-QE experiment. It tests a different hypothesis about
what information in the SOM may be useful for mesh adaptation.

\subsection{Interpretation of the refinement-utility analysis}

The cell-matched utility analysis provides a more nuanced interpretation
than simply comparing global mean and RMS values. In the common selected
region, Gradient produced lower mean post-refinement error in both the
\(1N\) and \(2N\) comparisons. Gradient also produced lower mean error in the
Gradient-only region.

In contrast, the SOM-interface selection showed lower mean error in the
SOM-only region for the \(2N\) case, and the SOM-interface solution also
showed lower mean error in the large region selected by neither method.
These observations indicate that the global improvement cannot be reduced
to a statement that every SOM-selected cell is individually more useful
than every Gradient-selected cell.

The result is consistent with the global coupling of the CFD equations.
Refinement changes the solution throughout the computational domain, so
the numerical value of a refinement decision must ultimately be evaluated
through its downstream effect on the complete solution.

The utility analysis therefore supports a more cautious interpretation:
the SOM-interface criterion appears to identify a different set of spatial
structures whose refinement can alter the global solution beneficially,
but the causal relationship between individual selected cells and global
error reduction remains unresolved.


\subsection{Engineering implications}

The study illustrates a broader methodological point for machine-learning
applications in computational engineering.

A machine-learning indicator can exhibit spatial structure, correlate with a
reference discrepancy, and identify physically meaningful flow regimes while
still failing to improve the final numerical solution.

For this reason, evaluation should proceed through a hierarchy of evidence:

\begin{enumerate}
    \item verify the underlying CFD solution;
    \item determine what information the machine-learning model actually
    represents;
    \item compare the proposed indicator against a conventional baseline at
    equal computational budgets;
    \item solve the adapted CFD cases rather than evaluating indicator
    correlation alone;
    \item assess the final quantities of interest and global error metrics;
    \item account for the complete computational cost before making
    efficiency claims.
\end{enumerate}

Within the original four-budget SOM-QE/Hybrid-v1 matrix, the conventional
Gradient indicator remains the most consistently accurate baseline. The
subsequent SOM-interface experiments provide a different and more promising
result in cell-matched tests, but they are not yet sufficient to establish
general superiority. The principal contribution of the combined analysis is
therefore the development and testing of a more precise hypothesis: learned
flow-regime information may be useful for adaptive refinement when it is
combined with physically meaningful measures of local variation and
evaluated through the resulting CFD solution.

